%% ma: L12-B14-0086
%% lop: 12
%% chuong: 5
%% bai: 14
%% dang: 12.14.36
%% loai: TLN
%% muc: VD
%% dap_an: "80"
%% nguon: "Ảnh giáo viên gửi 07/11/2026 (Đề thi thử THPT Hoàng Diệu & Sào Nam, Đà Nẵng - 2026), Câu 21"
%% kien_thuc: "Bài 14 (phương trình mặt phẳng chứa trục tọa độ), Bài 7 (hệ trục tọa độ), định lí Pythagore"
%% trang_thai: cho-duyet
%% ly_do: "Dạng 12.14.36 mới đề xuất, chờ giáo viên duyệt."
%% ghi_chu: "Tài liệu gốc không kèm đáp án; đáp án do Claude giải và kiểm bằng sympy. Đề gốc không nói rõ đoạn nối An với chỗ bóng rơi vuông góc với Ox (hình gốc chỉ vẽ), nếu không có điều kiện này thì không xác định được vị trí bóng rơi; đã thêm vào đề, và ghi rõ bóng bắt đầu bay từ điểm trên tia Oz như hình gốc. Hình vẽ lại bằng TikZ, bỏ ảnh người (trang trí). Đáp số 80 không phụ thuộc bóng lệch về phía y dương hay âm."
\begin{ex}
Trong giờ thể dục học về kỹ thuật chuyền bóng hơi, Bình và An tập chuyền bóng cho nhau. Ở một động tác Bình chuyền bóng cho An, quả bóng bay lên cao nhưng lại lệch sang bên trái của An và rơi xuống vị trí cách chỗ An đứng $0{,}5$ m và cách chỗ Bình $4{,}5$ m (đoạn thẳng nối chỗ An đứng với vị trí bóng rơi vuông góc với đường thẳng nối chỗ Bình và chỗ An). Chọn hệ trục tọa độ $Oxyz$ sao cho gốc tọa độ $O$ tại vị trí của Bình, vị trí của An nằm trên tia $Ox$, mặt phẳng $(Oxy)$ là mặt đất và quả bóng bắt đầu bay từ một điểm trên tia $Oz$, đơn vị trên mỗi trục là mét (tham khảo hình vẽ).
\begin{center}
\begin{tikzpicture}[x={(1.1cm,0cm)},y={(-0.9cm,-1.1cm)},z={(0cm,1cm)}]
  \draw[truc] (0,0,0)--(5.6,0,0) node[below right]{$x$};
  \draw[truc] (0,0,0)--(0,1.8,0) node[left]{$y$};
  \draw[truc] (0,0,0)--(0,0,3.6) node[left]{$z$};
  \coordinate (O) at (0,0,0); \coordinate (An) at ({2*sqrt(5)},0,0); \coordinate (L) at ({2*sqrt(5)},0.5,0); \coordinate (S) at (0,0,1.6);
  \draw (O)--(L) node[pos=.55,below,font=\small]{$4{,}5$ m};
  \draw (An)--(L) node[pos=.6,below right,font=\small]{$0{,}5$ m};
  \draw[khuat] (S) .. controls (1.4,0.15,3.6) and (3.4,0.38,3.2) .. (L);
  \fill[orange] (S) circle (2.5pt);
  \fill (An) circle (1.5pt) node[above]{An};
  \node[left] at (S) {Bình};
  \node[below left] at (O) {$O$};
\end{tikzpicture}
\end{center}
Biết rằng quỹ đạo của quả bóng nằm trong mặt phẳng $(\alpha)\colon x+by+cz+d=0$ và $(\alpha)$ vuông góc với mặt đất. Tính giá trị của $b^2+c+d$.
\shortans{80}
\loigiai{Vị trí bóng rơi $L\left(x_L;\pm0{,}5;0\right)$ với $x_L^2+0{,}5^2=4{,}5^2\Rightarrow x_L=2\sqrt5$.
$(\alpha)$ vuông góc với mặt đất và chứa điểm bóng bắt đầu bay trên $Oz$ nên chứa trục $Oz$: $c=0$, $d=0$.
$L\in(\alpha)\Rightarrow 2\sqrt5\pm0{,}5b=0\Rightarrow b^2=80$. Vậy $b^2+c+d=80$.}
\end{ex}
